\documentclass[../main.tex]{subfiles}
\begin{document}
%==========================================TIÊU ĐỀ TẬP========================================
\begin{table}[h]
  \begin{adjustwidth}{-1cm}{}
    \begin{tabular}{>{\centering\arraybackslash}p{6cm}|>{\raggedright\arraybackslash}p{12.5cm}}
      \multirow{5}{6cm}
      % Cột bên trái: ÉPISODE và số tập
      {\\[-40pt]
        \flaregothic\fontsize{20pt}{20pt}\selectfont \centering
        \color{eptype}PERSONNALISÉ\color{black}\\[12pt]
        \vnmdisp\fontsize{36pt}{36pt}\selectfont
        \color{epnum}SÁU
        \\[18pt]
        % \flaregothic\fontsize{14pt}{14pt}\selectfont
        % \color{epattr}BIENVENUE, \uline{\color{Banpire} Mel} !
        % \flaregothic\fontsize{18pt}{18pt}\selectfont
        % \color{epattr}ÉPISODE PERDU
        \color{black}
      }
       &
      % Dòng thứ nhất: tên tập tiếng Nhật
      {\fotskip\fontsize{18pt}{18pt}\selectfont さっき\ruby{廊}{ろう}\ruby{下}{か}で\ruby{見}{み}つけてきたの (\ruby{宇}{う}\ruby{宙}{ちゅう}276\ruby{号}{ごう})
      } \\[6pt]
      % Dòng thứ hai: tên tập tiếng Anh
       & {\cafeteria\fontsize{18pt}{18pt}\selectfont The Tale of Momotaro}                                                            \\[4pt]
      % Dòng thứ ba: tên tập tiếng Việt
       & {\vnmsans\fontsize{18pt}{18pt}\selectfont Cuộc phiêu lưu tại đảo Oni của Momotaro}                                           \\[8pt]
      % Dòng thứ tư: Nhân vật xuất hiện
       & {\sen Track used: Relics (\ruby{遺}{い}\ruby{跡}{せき})}
      \\[18pt]
      % % Dòng cuối cùng: Thời gian diễn ra của tập (thường sẽ là ngày phát hành)
      % & {\continuum\fontsize{16pt}{16pt}\selectfont Ngày phát hành: 12/09/2021}\\[18pt]
    \end{tabular}
  \end{adjustwidth}
\end{table}
%=========================================TRUYỆN CHÍNH========================================
\color{black}

\alt{128}{Mười}{276}

\begin{center}
  \uline{Trong tập này, A-chan là người kể chuyện, vì thế các câu được viết dưới đây đều được viết theo lời của cô.}
\end{center}

Ngày xửa ngày xưa, vào một buổi sáng se lạnh ở vùng quê yên bình thuộc thời Edo, có hai ông bà lão sống lặng lẽ bên bờ suối. Một ngày nọ, bà lão (vốn là Oozora Subaru, không hiểu vì lý do gì lại phải nhập vai một bà cụ) đang lom khom giặt áo kimono với tâm trạng mệt mỏi vì tuổi cao sức yếu\dots

``\hl{Tại sao em lại phải đóng bà già chứ? Em tưởng em là chim trĩ cơ mà\dots}'' bà già càu nhàu, hai tay vẫn đang chà xát bộ áo, khuôn mặt hiện rõ vẻ cam chịu.

Bà cụ phơi bộ kimono lên dây phơi, thở dài một hơi nhẹ, lẩm bẩm: ``\hl{Chăc là do kịch bản nên mới như vậy chứ\dots}''

Bỗng nhiên, từ thượng nguồn, một vật gì đó to và tròn trôi theo dòng nước. Một vật gì đó ửng đỏ hồng, mang dáng dấp của một thứ quả.

Một quả đào khổng lồ!

Subaru há hốc miệng nhìn, thốt lên: ``\textbf{Ể!? Quả đào to thế!?}''

Đúng lúc ấy, từ phía bụi lau sậy, một cô gái tóc hồng nhỏ nhắn hớt hải chạy ra, tiếp cận tới quả đào đang trôi ấy.

``Su-ba-ru-cha-sen-pai\~{} Đưa quả đào đó cho em na\~{}!! Em mún làm nức đào dầm!'' cô bé hí hửng nói.

``K-Khoan, Luna! Em không có vai đoạn này mà!? Sao em lại-!?''

Đúng lúc đó, cụ ông - tức Hikawa Kuon - lọm cọm từ trên núi trở về, tay xách rổ củi. Ông khàn giọng, nhìn về bà cụ nói: ``Bánh kê ơi, em phải đến tận đoạn sau mới xuất hiện cơ.''

``Ể\~{}? Nhưng mà Luna mún ăn quả đào cơ-nanora!!'' cô bé bí ẩn lăn ra ăn vạ, nước mắt lưng tròng.

``Chốc nữa anh lên rừng tìm một quả đào khác còn to hơn thế này nhiều!'' ông già cần mẫn nói, dù hơi lên giọng đôi chút.

Bỏ qua những lời thắc mắc lẩm bẩm của bà lão và lời trấn an từ cụ ông, Subaru cuối cùng cũng lăn trái đào về tới túp lều nhỏ, nơi ông lão đang nhóm bếp, vừa mới quay trở về sau cuộc mặc cả với cô bé bí ẩn kia.

``Kuon\~{} Em kiếm được quả đào to lắm này!'' bà lão gọi vang lên từ hiên nhà.

Ông lão tất bật về phòng, nhìn thấy quả đào căng và mọng nước, to cỡ một chiếc bàn. Ông tỏ ra hí hửng (dù trước đó ông đã nhìn thấy quả đào từ trước, nhưng vẫn phải đóng như phim): ``Ồ, trái đào trông căng thật đấy!''

``Thế này là đủ cho vợ chồng mình ăn cả tháng rồi!'' Subaru tiếp lời, hớn hở không kém gì ông lão.

Kuon vội chạy vào bếp, lấy ra một con dao đã cùn từ lâu (đạo cụ cho trước, theo yêu cầu của Kuon), chuẩn bị xẻ trái đào.

Nhưng vừa khi chiếc dao vừa khứa nhẹ lớp vỏ, một tiếng *rắc* vang lên như bị vỡ, khiến cho ông bà lão giật mình.

Quả đào mở ra, bên trong là dáng dấp một sinh linh vừa mới thức dậy sau giấc ngủ dài.

``U-Uaa\dots\ Sáng quá\dots\ Ai gọi em dậy vậy\dots''

Một bé gái xinh xắn với đôi tai sói mềm mại và cặp mắt vẫn còn lim dim buồn ngủ. Cô giương đôi mắt hổ phách to và tròn nhìn hai người.

``\dots Ể? Là con gái à?'' ông lão thốt lên đầy ngạc nhiên.

``Không phải Momotaro là con trai sao!?'' bà lão ngạc nhiên, tới mức ngã bật ra đất.

Đó chính là Mio, cô bé được sinh ra từ quả đào. Nhưng để tránh nhầm lẫn, ta hãy gọi cô bé là Mio bé con.

Cô bé Sói đen vẫn ôm quả đào như chăn, ngước nhìn mọi người: ``Em\dots\ ăn xáng chưa í nhỉ?''

Subaru đứng đơ người, nhìn về ông lão: ``Cái này\dots\ em thấy lo rồi đấy.''

Kuon nhìn tình cảnh hiện tại chỉ biết thở dài, tay chống nạnh: ``Thôi thì\dots\ đành nhận nuôi luôn đi vậy.''

Và từ hôm đó, ông bà lão hết lòng chăm sóc cô bé như con ruột. Mio lớn nhanh như thổi- ờm\dots\ ít nhất là về tinh thần. Chứ ngoại hình thì vẫn y như ngày cô bé chui ra từ trái đào: nhỏ nhắn, má phúng phính, tóc đen dài chấm vai, và đôi tai sói lúc nào cũng vểnh lên hóng chuyện.

Ban ngày, cô bé nghịch ngợm trèo cây, đào hang, nhảy múa trên ruộng như những đứa trẻ mới lớn khác ở trong làng. Có lẽ bạn sẽ không tin, nhưng trẻ con ở đó lại rất hòa hợp với cô bé, mặc cho việc cô bé không hoàn toàn là con người.

Nhưng khi màn đêm buông xuống, mọi chuyện\dots\ lại trở nên có phần hỗn loạn tới ngớ ngẩn hơn nhiều. Nhất là vào những đêm trăng tròn.

``\textbf{Awooooooooo\~{}\~{}\~{}\~{}\dots}''

Tiếng hú rất to vọng tới cả làng, tới mức dân làng hoảng loạn, vội khua hết tài sản về nơi an toàn. Một người dân trong làng thốt lên: ``\textbf{Quỷ đến! Quỷ về làng rồi bà con ơi!!}''

Nhưng cũng có những người nhận ra tiếng hú này. ``Bình tĩnh nào mọi người! Là con bé Sói con nhà ông bà Subaru-Kuon đấy! Nó lại hú!''

Mio bé con tiếp tục hú, tới mức cả ông bà lão không tài nào ngủ được. Bà lão bật dậy, bịt tai lại, hét lên như mắng: ``Trời ơi, Miosha! Khuya rồi mà hú thế là hàng xóm lại sang mắng đấy!''

Kuon ở bên cạnh, trằn trọc không ngủ được. Ông chỉ thở dài, như thể đã quá quen với chuyện này: ``Lúc đầu anh tưởng em ấy là loài đêm\dots\ Giờ thì chắc chắn luôn rồi.''

Dẫu cho việc những đêm phiền hà như thế, Mio bé con vẫn được cả dân làng yêu quý. Trong suốt nhiều năm sau đó, cô bé vừa sống như một đứa trẻ bình thường, vừa hú đêm như một thói quen. Cô bé vẫn giữ tính vui vẻ, ngây thơ và hoạt bát, ngày càng hòa đồng với mọi người. Mặc dù vẻ ngoài vẫn bé như chiếc gối ôm, tâm hồn Sói con đã cứng cáp hơn rất nhiều.

Một ngày nọ, tin dữ lan đến tai dân làng, làm đảo lộn hoàn toàn cuộc sống của ngôi làng yên bình này.

``\textbf{Quỷ từ đảo Onigashima lại xuất hiện rồi! Chúng cướp gạo, đốt làng, bắt cả trâu bò!}''

Bộ tộc Quỷ từ hòn đảo kế cận đã quay trở lại. Ngày càng có nhiều hung tin nhận tới hơn: có ngôi làng ở cạnh bên thì thiêu rụi hoàn toàn trong biển lửa, có ngôi làng thì bết bát mùa thu hoạch do bị bọn Quỷ cướp sạch, có gia đình thậm chí còn bị bắt mất trâu bò, tài sản mà cả đời họ đã tích cóp được.

Bên trong gian nhà, vợ chồng Subaru-Kuon đang ngồi nói chuyện, tâm lý hoàn toàn bị động trước hung tin trên.

``Không lẽ lần này tụi nó định sang tận đây?'' bà lão rụng rời tay chân, mặt tái mét vì sợ hãi.

``Dễ là như vậy lắm\dots'' ông lão tỏ vẻ cẩn trọng. ``Đợi đó, để anh chuẩn bị đồ-''

Bất chợt, Mio bé con nắm chặt tay, tuyên bố như một vị anh hùng (đúng theo cả nghĩa đen và nghĩa bóng): ``Nếu như thế\dots\ Em sẽ đi đánh Quỷ!''

``\dots Hả?'' cả hai ông bà ngước nhìn ``bình rượu mơ'' của mình, há hốc mồm kinh ngạc.

``Khoan! Em muốn nói rằng là\dots\ em muốn đi một mình ư!?'' bà lão hỏi lại, mắt vẫn trố ra.

Cô bé Sói đen gật đầu, đôi mắt hiện lên một ngọn lửa quyết tâm. Dù ông bà lão có khuyên thế nào, cô bé vẫn kiên quyết. Cô ấy chạy ra giữa làng, đứng trên một tảng đá, tay chỉ thẳng về phía chân trời, hét lớn gây chú ý cả dân làng: ``\textbf{Cháu sẽ đi đánh Quỷ! Không thể để tụi nó làm hại thêm ai nữa!}''

Một người đàn ông dân làng chỉ về phía bé Sói, rồi chỉ về cặp vợ chồng vừa chạy ra khỏi gian lều để đuổi theo: ``Ô-Ông bà muốn để một con vịt giời á nhân này đi xung trận ư!?''

``Cô bé này có làm được trò trống gì không vậy\dots?'' một người phụ nữ khác lên tiếng, tỏ vẻ nghi ngại. ``Nó lại là con gái nữa chứ\dots''

Mio bé con phồng má lên, tỏ vẻ giận hờn: ``Cháu là Sói! Cháu thực sự có thể hú đuổi chúng đi!''

Trong lúc bé Sói và dân làng vẫn đang cù cưa với nhau, ông lão thì thầm với bà lão: ``\hl{Anh không nghĩ là đám Quỷ sợ tiếng hú đâu.}''

Bà lão đổ mồ hôi hột: ``Nhưng chắc là chúng sợ\dots\ tiếng hú lúc giờ Sửu\footnote{Từ một giờ sáng tới ba giờ sáng.}.''

Vậy mà không ai ngờ với lòng quyết tâm sáng như trăng rằm và tiếng hú vang động cả làng mỗi đêm, Mio bé con cuối cùng đã thuyết phục được mọi người ủng hộ cô bé một cách dễ dàng. Người thì cho bánh, người thì cho kiếm gỗ, người thì dúi cho con cá khô với mảnh bùa.

Tất cả mọi người trong dân làng đều chúc phúc cho cô bé Sói đen tuy nhỏ tuổi nhưng dũng cảm này.

Và thế là, cuộc hành trình của ``Momotaro phiên bản Sói con'' bắt đầu, mang theo niềm hy vọng của cả dân làng.

\ct

Từ ngày rời làng, Sói con Mio đã đi mãi, đi mãi, qua bao nhiêu là ruộng bậc thang, rừng tre rậm rạp và cả một đàn bò đang\dots\ ngồi thiền. Cô mang theo ước mơ bảo vệ người dân và một túi đựng bánh kê, thứ bánh được làm bởi một cô bé bí ẩn, nghe đồn là ở thị trấn làng bên.

Hiển nhiên, một cuộc hành trình để đánh bại lũ Quỷ gặp rất nhiều gian truân. Một trong số đó là việc cô bé cần phải có những người sát cánh yểm trợ cho cô.

Chính vì thế, cô bé Sói đen đã tự nghĩ thầm với bản thân: ``Phải tìm được những người bạn trước đã thì mới có thể đánh quái!''

Với chiếc kiếm gỗ cầm trên tay và một túi đựng bánh kê và tấm bản đồ vẽ bằng sáp màu, em ấy hừng hực khí thế rảo bước giữa con đường đất gồ ghề dẫn ra biển lớn, nơi mà bọ tộc Quỷ hung ác đang trú ngụ\dots

``\textit{Không biết liệu cô bé có kết nạp được đồng đội nào không, hay lần này đi kiểu một cân ôn đây?}'' Kuon giờ đã hết vai, ngồi ở băng ghế khán giả.

Cô bé Sói đen rảo bước trên cung đường đất, ngân nga giai điệu đồng dao của trẻ con. Cô dừng chân lại khi thấy một cái xe đẩy bánh nhỏ được dựng bên đường.

``Ồ? Xe bánh này của ai ta?''

Mio bé con đi vòng ra phía trước, bắt gặp một cô bé lùn lùn, tóc hồng được tết hai bên, mặc đồ công chúa nhưng đang đội khăn bếp. Trên xe là một loạt bánh tròn bé tí, bốc khói thơm lừng.

``Oa\dots\ Thơm quá! Em đang bán gì đó?''

Chủ xe tươi cười đáp: ``Bánh kê đó-nanora! Bánh do em nướn bàng đôi tay tràn đầy sự iu thưn này!''

``\textit{Không biết tại sao bánh kê lại có thể hóa thân thành một cô công chúa nhí. Có lẽ do ``mẹ kể chuyện'' (A-chan) viết, hoặc là do ông Cường chuyển thể rồi.}''

``Vậy\dots\ em là ai?'' Mio bé con nghiêng đầu.

``Em là Kibi dango xống! Chị có thể gọi em là `Dango-sama' cũng được na! Còn chị là\dots?''

``Em là Momotaro! À không, là Sói Taro thì đúng hơn!''

``\textit{Tên chính thức là Momotaro, nhưng có vẻ em ấy muốn gọi với biệt danh hợp với loài vật hơn,}'' Kuon tiếp tục bình luận.

Cô chủ xe cười tít mắt. ``Nghe hay đó-nora. Zậy em có bánh kê, em mún mời chị một cái! Ăn vào là có xức mạn đấy-nanora!''

Đôi mắt háu ân của Mio bé con sáng rực lên. ``Thật á!? Vậy cho em xin một cái!''

Luna tung chiếc bánh kê của mình lên cao, như thể là cách mà người ta huấn luyện chó. Cô bé Sói đen nhảy lên cao, ngoạm lấy miếng bánh đầy ngon lành với nụ cười nở trên môi.

Ngay khi cô bé vừa nuốt xong miếng đầu tiên, một tiếng hú dài cất lên từ rừng rậm bên cạnh. Một tiếng ``Awwoooooo\~{}!!'' vang vọng cả khu đồi, làm đôi tai nhỏ nhắn của Sói Taro khẽ vểnh lên như gặp được đồng loại.

``À-rế!? Tiếng gì vậy?'' bé Sói thốt lên, mắt ngước về cánh rừng.

Một bóng người bước ra khỏi rừng, đó là một cô gái lớn hơn Mio-loli, mái tóc đen tương tự, tai sói vểnh lên, nụ cười hiền hòa. Nói không ngoa, đây là ``phiên bản trưởng thành'' của bé Sói con.

Hai cô nàng - một bé, một lớn - nói chuyện với nhau qua những tiếng hú của mình. Để tiện theo dõi hơn, họ sẽ nói chuyện với nhau bằng tiếng con người.

``Em đang định đi đánh Quỷ à?'' Mio lớn hỏi, bản mặt không có ý dọa dẫm gì, mà trái lại rất hiền hòa.

Mio bé con há hốc mồm: ``C-Chị là ai vậy!? Sao nhìn giống em thế!?''

``Chị là Sói như em thôi. Nhưng lớn hơn một chút. Có thể gọi chị là Chó cũng được, nếu theo truyền thống.''

Cô bé đẩy xe bánh ngơ ngác nhìn hai nàng Sói hú nhau, không hiểu họ đang nói gì. ``Ơ-Ờm\dots\ Taro-chan? Nó đang nói gì thế-nora? Chị khum hỉu gì hớt á\dots''

``Chị có muốn cùng chúng em đi đánh Quỷ không?'' Mio bé con hỏi, giọng có đôi phần háo hức, ngó lơ hoàn toàn câu hỏi của Luna.

``Nếu em cho chị một cái bánh kê. Có vẻ như bạn của em còn đang ngơ ngác không hiểu chúng ta đang nói gì kìa.''

Cô bé Sói nhìn Chủ xe, rồi chia cho Sói lớn (hoặc Chó, nếu diễn theo đúng vai) một chiếc bánh, mong rằng mọi người có thể hiểu nhau hơn.

Mio lớn vẫy đuôi quần quật, khen nức nở: ``Bánh kê ngon thật\~{} À, chị nói được tiếng người rồi này!''

Bé Sói sau đó kể lại những gì mà hai người họ vừa trao đổi, và Luna cũng gật đầu đồng ý. Thành viên đầu tiên gia nhập cuộc hành trình, một người bạn đáng tin cậy với cú hú gây náo loạn cả vùng quê vào buổi đêm, chính là chị Chó!

``\textit{Chính xác hơn là Sói đen,}'' cậu bình luận, ``\textit{nhưng mà thôi kệ, thay đổi nguyên tác một chút chắc cũng không sao nhỉ?}''

\ct

Cả ba người tiếp tục lên đường. Không lâu sau, họ thấy một cô gái tóc xanh đang đu dây từ cành cây này sang cành khác, miệng lẩm bẩm gì đó khó hiểu.

``Ơ? Ai kia?'' Miotaro chỉ về người đó.

Cô gái nhìn giống khỉ này treo ngược lên cành cây, hú khẹc đầy hứng khú. Nhìn thấy chiếc xe bánh mà Luna đang đẩy, đôi mắt của cô nàng Khỉ bỗng chốc sáng lên.

``Cô ấy chắc cũng bị hấp dẫn bởi bánh kê thần kỳ này đấy,'' Chó (sói) Mio lên tiếng, đọc vị được cô nàng ở trên cây.

Momotaro gật đầu ra vẻ hiểu. Cô lấy ra một chiếc bánh trên xe, đưa nó trước mặt của Khỉ, hỏi với giọng chân thành: ``Chị có muốn một cái không ạ?''

Khỉ tóc xanh vui vẻ, với lấy chiếc bánh mà Mio bé con đưa cho, tới mức suýt trượt tay khỏi cành cây. Cô cho nguyên cái bánh vào miệng, nhai đầy ngon lành.

Bất chợt, đôi mắt sao băng của cô bừng sáng. ``Whoa! Chị\dots\ chị nói được tiếng người rồi sao!?'' Suisei bất ngờ với bản thân khi cô cuối cùng có thể giao tiếp như một con người.

``Hịu quả dõ dệt luôn-nanora! Em còn nhìu lắm nha!'' Luna hớn hở, đưa cho cô Khỉ thêm một chiếc bánh nữa.

Trong lúc Suisei ăn chiếc bánh thứ hai với vẻ mãn nguyện, Mio bé con hào hứng hỏi: ``Chị Khỉ ơi! Chị có muốn đi cùng chúng em không? Chúng em đang trên đường đi đánh Quỷ đấy!''

``Đánh nhau à!?'' cô Khỉ hơi cao giọng, pha chút trong đó một sự phấn khích. ``Lâu lắm rồi mới có cớ để luyện giọng đấy!''

``Nếu chị thích thì có thể gia nhập với bọn em! Đảm bảo là vui lắm đó!'' bé Sói nói, tiếp tục kích thích cô.

``Được rồi, lên đường nào!'' Suisei nhảy xuống mặt đất, giơ nắm đấm lên cao, khí thế hừng hực. ``Chị sẽ cho mọi người thấy Hoshimachi Suisei này ghê gớm như thế nào!''

\ct

Họ tiếp tục cuộc hành trình tới đảo Onigashima, với bốn thành viên đang dạo bước, trò chuyện với nhau đầy vui vẻ. Họ tiếp cận đến một bãi đất trống.

``Shuba-shuba-shuba!!'' một tiếng kêu vang lên ở một trong những bụi cây quanh đó.

``Tiếng này là\dots\ vịt ư?'' nàng Chó (sói) nghiêng đầu, cảm thấy ngờ ngợ vì tiếng kêu quen thuộc.

Một con chim trĩ khoác trên mình bộ lông xanh, chiếc mỏ (gắn bằng đồ hóa trang) màu đỏ đang kêu lên một thứ tiếng kỳ quặc mà một đồng loại sẽ không bao giờ kêu.

``Chị ấy\dots có phải là chim trĩ khum zậy-nora?'' Luna đổ mồ hôi hột.

``Hay là vịt trời?'' Suisei chỉ tay về phía con chim trĩ vẫn đang kêu.

``Hoặc là\dots\ ngan?'' tới lượt Mio bé con đoán.

``\textit{Cái vụ này thì anh không chắc là chim trĩ kêu như thế nào nữa,}'' Kuon ở bên cánh gà nhún vai, ``\textit{nhưng thôi kệ, đều là chim cả.}''

Cô chim trĩ vẫn đang kêu lên những tiếng ``shuba shuba'' đầy kỳ quái, nhưng càng lúc càng lớn hơn, như đang để cố thanh minh rằng: ``Chị không phải là vịt! Chị là chim trĩ!!''

``Có khi cho chị í cái bánh đi, biết đâu chị í nói chuỵn đực á-nora?''

``Ý hay đó, Chủ xe!''

Luna đưa một chiếc bánh kê khác cho Mio bé con, rồi cô bé đưa bánh cho cô chim trĩ. Không một chút do dự, cô chim trĩ đã cắp chiếc bánh rồi ăn ngấu nghiến như thể chưa ăn nhiều ngày.

``Shuba- À-rế!? Chị nói được tiếng người rồi á!? Không phải `shuba shuba' nữa à!?'' Subaru - tên của cô chim trĩ lai vịt này - thảng thốt khi cô có thể nói được tiếng người.

Khỉ Suisei đập tay với Chủ xe: ``Tuyệt vời! Thêm một thành viên nữa! Đội mình giờ bắt đầu đông vui rồi!''

``Ủ-Ủa!? Chị còn chưa hiểu chuyện gì đang xảy ra mà!''

``Thôi nào Subaru, bọn chị biết là em sẽ đồng ý tham gia cuộc hành quân tới đảo Onigashima rồi,'' Chó Mio cười hiền, như thể cô đoán được cô sẽ định hỏi gì.

Cuối cùng, như thể bị ``người mẹ dịu hiền'' bắt bài, chim trĩ Subaru đành phải đồng ý gia nhập vào nhóm, khiến cho Mio bé con vừa cảm thấy an tâm hơn lại vừa vui hơn vì có thêm bạn mới.

\ct

Và thế là, từ một cô bé Sói con đơn độc, nhóm bạn giờ đã trở thành một đội hình đầy đủ, một Sói, một Chó, một Khỉ, một Chim trĩ\dots\ và một ``chiếc bánh'' mang hình hài của một công chúa, nói cười đầy nũng nịu.

``\textbf{Chúng ta đến bừ biển rồi!}'' Miotaro giơ hai tay lên cao, reo lên phấn khởi.

``Nhưng vấn đề là\dots'' Chó Mio lập tức nêu ra vấn đề trước mắt.

Họ chỉ còn một chướng ngại vật nữa trước khi thực sự đặt chân vào thử thách thực sự: đại dương bao la chia cắt giữa đất liền và nơi ẩn náu của bộ tộc Quỷ.

Đứng trước bờ biển với bãi cát vàng óng từ nắng chiếu rọi xuống và tiếng sóng vỗ dạt dào, cả đội đều đứng trước một bài toán hóc búa: làm thế nào để có thể băng qua bên kia.

``Ta phải bơi hả?'' Suisei gãi đầu gãi tai.

``Chị không biết bơi đâu nha!'' Subaru thốt lên. ``Chị là chim, không phải cá!''

``Nhưng mà chị bay được mà, Chim trĩ-onee-chan?'' Mio bé con nghiêng đầu.

``Chị không đủ sức để kéo mọi người sang bên kia.''

``Vậy chắc chúng ta cần một chiếc thuyền đấy\dots'' Mio lớn đề xuất.

``Để em thử nàm một thuỳn từ bánh kê nha-nanora?'' Luna hớn hở giơ tay, nhưng lập tức bị Chim trĩ lắc đầu: ``Nếu mà thế thì chưa kịp chèo thuyền đã rã ra rồi đó, Luna ạ.''

\il

``Việc gì cũng đến tay mình, tài thật\dots\ Hết đóng ông lão, giờ lại đóng rùa\dots\ Rùa đấy,''  cậu Tổng từ phía cánh gà càu nhàu.

``Anh là người duy nhất còn thừa vai mà, Kuon-senpai,'' Fubuki đang hóa trang trong bộ đồ Quỷ nhoẻn miệng cười.

``Với lại rùa còn có mai nữa. Còn ai cứng đầu ngoài cậu đâu?'' A-chan nói, cố giữ bình tĩnh nhưng cũng không thể nhịn cười.

Kuon nghe hai người nói thế, chỉ có thể thở dài bất lực: ``Được rồi, làm thì làm.''

\il

Trong lúc mọi người vẫn còn đang vắt óc suy nghĩ để xâm nhập được đại bản doanh của tộc Quỷ, thì từ giữa đại dương rộng lớn, hàng tiếng trống *TÙNG TÙNG* vang lên.

``T-Tiếng zì thế-nanora!?'' Chủ xe bánh kê giật mình.

``N-Nhìn kìa mọi người!'' Subaru chỉ cánh về phía vật thể đang nhấp nhổm giữa biển.

Trong làn nước xanh kia, một cái mai rùa khổng lồ  nhô lên khỏi mặt nước, dần dạt về phía năm cô gái. Cái đầu của con rùa trồi lên, hiện lên ánh mắt\dots\ chán chường không thể tả vì bị thúc ép, thay vì là một sự gì đó gọi là hiền từ và muốn giúp đỡ.

Một gương mặt thân quen với bé Sói Momotaro của chúng ta.

Con rùa ấy nhìn năm cô gái, nói với giọng bằng bằng: ``Xin giới thiệu: tôi là con rùa từ truyện khác, bị kéo qua làm cameo cho vở này. Ai thấy tôi sang thì nhớ mời tôi bánh kẹo sau nha.''

Miotaro lóa sáng mắt lên, vỗ tay đầy phấn khởi: ``Oa\~{}! Một con rùa khổng lồ kìa! Bác Rùa ơi, bác đưa chúng cháu sang bên kia được không ạ?''

Chó (sói) Mio nghiêng đầu: ``Rùa\dots\ biết nói ư? Không cần bánh dango luôn à?''

Rùa Kuon thở dài đầy mệt mỏi, nói như đọc thoại: ``Ừ, biết nói, biết bơi, và biết cáu nữa. Lên hết đi để còn đẩy nhanh tiến độ.''

Tất cả mọi người, không ai nói ai, lập tức leo lên mai rùa, ngoại trừ Luna, người vẫn còn lưỡng lự. ``Zậy\dots\ chở lun xe bánh kê của tui đực hơm?''

``Chở luôn. Nhưng mà nếu anh chìm thì đừng có kêu nhé.''

``Kuoncha-senpai khéo lo-nora! Xe này có dây buộc mà!''

Thế là, nhóm được một con rùa (dù khá bất đắc dĩ) giúp sức, đưa cả năm người họ cùng chiếc xe vượt qua đại dương, hướng về phía đảo Onigashima. Chiếc xe bánh kê của Luna thỉnh thoảng cứ trượt đi trượt lại như một món hàng đang bấp bênh.

``\textit{Chẳng hiểu sao mình lại ở đây\dots}'' Kuon nghĩ thầm, cố gắng kìm nén sự mệt mỏi vì mọi thứ đã bắt đầu manh nha sự lộn xộn, dù vẫn ở trong tầm kiểm soát.

\ct

Năm chiến binh tí hon đã bộ lên bãi đá, nơi mà bộ tộc Quỷ đang ẩn náu, tiếp tục lên kế hoạch những chiêu trò tàn độc khác nhằm gieo rắc nỗi sợ cho dân chúng.

Sứ mệnh của nhóm không gì hơn ngoài việc chấm dứt những hành động hung ác này.

Nhóm anh hùng nhìn thấy ba cô ác quỷ đang tụ họp giữa đống kho báu lung linh, cùng với hàng tấn lúa mì và trâu bò mà chúng cướp bóng dược.

Một con quỷ tai cáo với chiếc kính râm và máy tính tay (cả hai thứ này đều không thuộc về thời đại ở thế kỷ XVIII), cười ranh mãnh với đống chiến lợi phẩm thu được: ``Chúng ta cần tối ưu hóa mô hình bán hàng! Cướp đồ rồi bán lại. Kiểu gì cũng lời gấp ba lần luôn đó!''

``\textit{\dots Bằng một cách nào đó, chủ nghĩa tư bản lại tồn tại được ở thời này\dots\ hình như câu chuyện này được sáng tác trước cả cách mạng công nghiệp lần thứ nhất (1784) nhỉ?}'' Kuon lúc này đang thu mình trong mai rùa, lẩm bẩm trong lúc đợi nhóm anh hùng.

Ở bên cạnh Quỷ Cáo trắng là một cô Quỷ khác nhỏ con hơn, nhưng đôi mắt cô đã ngấn lệ từ khi nào. Cô ôm lấy chiếc gối từ đống đồ đó, thỉnh thoảng sụt sịt: ``Không một ai khen em dễ thương cả\dots\ Tại sao mọi người lại xa lánh chứ\dots\ Hức\dots''

``Tôi thấy bà vẫn xinh đẹp mà,'' một con Quỷ tai nhọn ở bên cạnh, vỗ về như một người mẹ dỗ dành đứa trẻ. ``Nếu bà ngừng khóc một chút thì sẽ kiểm soát cảm xúc của mình tốt hơn đấy.''

Từ sau mỏm đá, Mio bé con và nhóm bạn đang nhìn ba ả Quỷ ở đằng xa. ``Vậy đó là ba kẻ đầu sỏ à\dots''

``Kế hoạch là thế nào đây, Miotaro-chan?'' Suisei hỏi khẽ, mắt đảo đi đảo lại giữa nhóm bạn và lũ quỷ.

``Em cùng với chị Chó sẽ cùng nhau hú để hù dọa bọn họ! Em tin là bọn chúng sẽ sợ thôi!'' Momotaro đáp.

``Em có chắc là được không vậy?'' Chim trĩ nghiêng đầu, tỏ vẻ lo ngại.

``Đảm bảo là được!''

Thế rồi, bé Sói hắng giọng, rồi bắt đầu cất tiếng hú lên như gặp phải trăng rằm. Chó (sói) Mio theo bản năng cũng giương mặt lên trời cất tiếng, tới mức những người còn lại bịt tai lại.

Tiếng hú của cặp đôi Sói hóa ra lại rất hữu dụng, khi nó động tới một trong ba con Quỷ ở phía bên kia, đặc biệt là Quỷ mít ướt. ``C-C-Có phải là tiếng của Quỷ không-nanodesu!? Đó chắc chắn là quỷ! Em tưởng có mỗi chúng ta là quỷ chứ!?'' cô nàng xanh mặt, ôm chặt chiếc gối của mình hơn.

``\dots Tôi dám chắc đó là tiếng hú của sói. Không phải một con đâu, mà tới tận hai con đấy,'' Quỷ tri thức dỏng tai lên nói.

``Bảo sao mà chị thấy tần số kỳ quặc thế,'' Quỷ tư bản rướn mày. ``Như thể đây là sóng siêu âm vậy.''

Nhân lúc cả ba con quỷ đang bị động bởi tiếng hú, cả năm cô gái nhảy khỏi mỏm đá, chạy lên tiếp cận ba tên đầu sỏ, với cặp đôi sói - một bé, một lớn - lên dẫn đầu.

Tiếng hú của hai người lớn tới mức người ta còn có thể nhìn thấy sóng siêu âm, làm nứt vỡ một phần của hang. Quỷ mít ướt bỏ chiếc gối, hét lên vì phát điên: ``'\textbf{Ồn quá! Ồn quá!!! Tha cho tôi đi mà!!!}'' Cuối cùng, con Quỷ đen đủi kia chạy vội vào hang, trên đường nó còn trượt ngã vài lần.

Giờ chỉ còn lại hai con quỷ. Luna đứng ở đằng sau\dots\ xem mọi người xung trận, không có mong muốn động chân động tay với họ. ``Các chị cứ đánh nhau thựt mạnh vào-nanora!''

``\textit{Mình không biết liệu Luna-chan muốn tạo cử chỉ hòa bình hay là thực sự không thể đánh đấm được nữa\dots}'' Rùa Kuon vẫn ở bên ngoài, đầu thò khỏi mai xem tình hình.

Ở sườn phía Tây của hang động, trận chiến giữa con Quỷ tri thức và con Chim trĩ bắt đầu nổi lên. Một bên là sử dụng kiến thức, và một bên là hỗn loạn, hai bên họ như lửa với nước.

``Ngươi là kẻ tsukkomi (thanh niên nghiêm túc) mới hả?'' Flare - cũng là tên !uỷ tri thức đó - trừng mắt nhìn Subaru. ``Ngươi trông có vẻ lòe loẹt nhỉ?''

``Thì ta là chim trĩ mà,'' Subaru nói. ``Ta đến đây để dạy ngươi biết thế nào là ngắt pha đúng lúc!''

Hai thanh niên nghiêm túc bắt đầu ``khua võ mồm'', gào thét tới mức ``ra lửa'' trước mặt đối phương. Thế rồi, hét nhau chán chê, họ bắt đầu va vào nhau như sấm sét, với một bên vung sách (cướp được từ lũ trẻ của dân làng) và một bên là vung cánh (bộ trang phục), đánh nhau tán loạn như lũ trẻ con mẫu giáo. Cả hai bên chạy vòng vòng, thỉnh thoảng vung nhau tới mức ``sứt đầu mẻ trán'' theo nghĩa bóng. Mỗi lần va chạm là một hiệu ứng *PANG!* như trong phim kiếm hiệp rẻ tiền.

``\textbf{Lập luận logic là vũ khí tối thượng!}''

``\textbf{Không phải lúc nào cũng sách vở như thế! Nhiều khi phải mơ tưởng một chút!!}''

*BỤP!*

Subaru tung một cú đá trời giáng vào thân khiến cho Flare lăn vào đống sách vở, mắt xoắn vòng tròn vì bị choáng.

\ct

Ở góc phía Đông của sân, một ca sĩ lừng danh của chi Vượn đang đối đầu với một\dots\ nhà buôn Cáo trắng lắm chiêu.

Fubuki - kẻ đầu sỏ thật sự của cứ điểm này - xòe ra\dots\ những bản hợp đồng và tập tài liệu để khè đối thủ: ``Ta có hơn mười hợp đồng thương mại, bốn mạng lưới phân phối và có lượng tài sản kếch xù! Nhà ngươi có gì nào?''

Suisei hất đôi tóc xanh của mình lên, giơ chiếc míc-rô mà cô đã chuẩn bị sẵn lên, nói với tông giọng của một kẻ máu lạnh: ``Ta có âm nhạc.''

Ngay lập tức, hai thùng loa bất chợt xuất hiện, cùng với hệ thống ánh đèn sân khấu sáng rọi lòe loẹt chiếu xuống, biến đại bản doanh của tộc Quỷ thành một cái sân khấu mini của riêng mình.

``Nhưng ta có hợp đồng nào với âm nhạc đâu?''

``Vậy thì ta sẽ cho ngươi sáng mắt ra.'' Nàng Khỉ xoay chiếc míc-rô của mình lên, rồi bắt đầu cất giọng ca của mình.

Âm thanh của bài hát vang lên dồn dập, làm rung chuyển toàn bộ hòn đảo, làm cả chấn động bên trong hang, nơi mà Quỷ mít ướt vẫn đang co ro bên trong.

Giọng ca trời giáng của Khỉ Suisei khiến cho tất cả mọi người - cả phe ta và phe địch - xao lòng. Không nói không rằng, toàn bộ các tộc Quỷ đều xán lại gần và\dots\ cùng chung hòa nhịp với cô nàng có đôi mắt Sao băng ấy.

``C-Các ngươi dám!? Làm thế nào mà các ngươi có thể bị giọng hát đó thu phục được\dots?'' Fubuki sốc nặng, hoàn toàn cảm tưởng như bị phản bội.

Khỉ không trả lời, tiếp tục hớp hồn mọi người bằng giọng ca của mình, như thể cô nàng đã trông chờ vào thời điểm này từ rất lâu.

Và rồi\dots

*LÓE!!!*

Ánh sáng đa màu rọi thẳng vào đôi mắt Cáo trắng khiến cho cô ngã ngửa ra đất. Thêm một con quỷ đầu sỏ nữa bị hạ gục bởi sức mạnh âm nhạc. ``Chắc\dots\ chắc ta nên đầu tư vào ngành giải trí\dots'' cô nằm lăn lóc giữa kho báu sáng choang, gồm bát sứ, kho báu, bảo vật, trâu bò và cả hàng tấn lúa gạo\dots

Nhìn thấy kẻ cầm đầu đã nốc-ao, Quỷ mít ướt lóc cóc từ hang chui ra, giơ hai tay lên: ``Được rồi\dots\ Bọn tôi đầu hàng\dots\ Làm ơn đừng hú nữa\dots''

Quỷ tri thức cũng lồm cồm bò dậy, xoa trán như vừa tỉnh dậy: ``Bọn chị thua rồi, Momotaro\dots\ À không, là Miotaro mới đúng.''

Như vậy, cả ba con Quỷ hùng mạnh nhất đã hoàn toàn bị nhóm anh hùng hạ gục. Chúng đã chính thức tuyên bố đầu hàng, dưới sự chứng kiến của năm chiến binh nhỏ nhưng đầy hỗn loạn.

\ct

Nhưng, niềm vui của người này lại là nỗi khổ của người khác. Không, không phải tộc Quỷ, chúng đã bị đánh bại tâm phục khẩu phục rồi.

Mà chính là con Rùa tên Kuon, nhân vật chở nhóm anh hùng ra đảo và đáng ra không nên xuất hiện trong câu chuyện này. Nhìn hàng chục túi báu vật được buộc thành chồng cao quá đầu, cậu há hốc mồm tới mức ngán ngẩm: ``Mình anh chở không hết được đâu\dots\ Đám này ăn cướp nguyên cả mấy làng xung quanh cơ à!?''

Chim trĩ chống hông, nét mặt giả vờ tỏ căng thẳng, nhưng không giấu nổi lo lắng: ``Thế anh định làm gì? Quay lại hơn mười chuyến nữa ư? Cơ mà anh ấy nói cũng có lý đó.''

``Anh là rùa, không phải tàu chở hàng đâu,'' cậu Rùa đáp. ``Có khi ai đó phải nhờ mấy người bên đất liền vác về hết thôi.''

``Thì Miotaro-chan đang quay zề đất lìn đấy thui-nora,'' Chủ xe bánh kê nói. ``Hơm biết chị í thế nào rồi nhỉ?''

Luna vừa dứt câu, một giọng nói cao vút vang lên từ phía xa: ``\textbf{Mọi người ơi!! Em quay về rồi đây!}''

Cô bé không hề đi một mình, mà đã nhanh nhạy nhờ các ngư dẫn vạm vỡ tới để đưa của cải của họ trở về nhà. Nét mặt ai cũng rạng rỡ khi cuối cùng họ đã lấy lại những gì của mình.

``Ủa? Ủa? Ủa? Mấy ông này ở đâu ra vậy?'' Suisei nghiêng đầu, nhìn những người đàn ông to con cập bến.

``Em mượn mấy bác ngư dân dưới làng đó ạ!'' Mio bé con hớn hở nói. ``Em có kể rằng em đã đánh bại Quỷ, thế là mọi người cùng xung phong ra khơi đấy ạ!''

Rùa gãi đầu gãi tai, cố đưa ra một lập luận của mình: ``\dots Vậy sao lúc đầu em không nhờ họ chở ra đảo luôn?''

Cô bé Sói nhỏ cười khì: ``Tại hồi đó em chưa nghĩ ra\dots\ Nhưng mà rùa dễ thương hơn mà.''

Kuon ngồi phịch xuống, than thở: ``Ôi, đời rùa đúng là khổ thật đấy\dots'' Tất cả mọi người đều phá lên cười.

Những tiếng hò reo, những tiếng hô lớn khi từng người chất túi hàng lên thuyền, khiến cho bầu không khí căng thẳng trên đảo Quỷ hồi nãy dần tan biến.

Không lâu sau, toàn bộ kho báu được chở về đất liền. Dân làng vui mừng khôn xiết, người người dựng lại nhà, tất cả mọi người càng quý trọng hơn vị anh hùng tài ba khi họ tôn vinh cô khắp mọi nơi, từ vòng hoa đặt ở nhà Subaru-Kuon, co tới biểu ngữ giăng khắp làng. Cuộc sống bình yên và êm đềm của ngôi làng đã dần được khôi phục.

Tất nhiên, những người dân trên đất liền không phải là những người duy nhất có được một cái kết viên mãn\dots

Tại hòn đảo Onigashima, nơi mà bộ tộc Quỷ đang sinh sống, các Oni ngồi bàn hội nghị với Luna - người xung phong ở lại hòn đảo - và Fubuki - trưởng tộc - về kế hoạch tái cấu trúc hòn đảo này.

Bên trong phòng họp được cải biên lại hoàn toàn, các thành viên trong tộc và cô bé chủ xe bàn tán với nhau rất sôi nổi về cách để thu hút du khách tới thăm.

``Xiên bánh kê mĩn phí cho mừi ngừi đầu tiên-nanora!'' Luna đề xuất.

``Chẳng phải là hơi ít sao?'' Flare chen vào, giọng ái ngại. ``Nếu muốn thu hút người ta thì ít nhất phải là một trăm xiên như thế.''

``Tôi có thể làm linh vật được không-nanodesu?'' Rushia chêm vào, háo hức khi cuối cùng được mọi người công nhận.

``Cái này được đấy. Dưới sự giám sát điều tiết hợp lý.''

Ở vị trí trưởng tộc, Quỷ tai cáo thì cười mãi không thôi về tham vọng làm giàu của mình: ``Chúng ta sẽ trở thành thiên đường du lịch Nhật Bản sớm thôi!''

``Ít nhất cũng phải dò xét tình hình xung quanh, rồi còn phải chờ người ta đóng dấu chấp thuận nữa chứ\dots'' Quỷ tri thức hỏi vặn lại.

``Ể\~{} Nhưng mà ta là Quỷ mà! Quỷ thì làm gì mà chả được!''

``Vậy nà chúng ta hợp tác với nhau thui nhỉ-nora?''

Và như vậy, hai bên đã bắt tay nhau, bắt đầu công cuộc xây dựng và tái thiết hòn đảo Quỷ này.

Từ phía xa bàn tròn, Chim trĩ Subaru trố mắt ngạc nhiên: ``Khoan đã! Chị tưởng mình đánh đuổi chúng chứ?''

Dường như Cáo trắng đã nghe được, liền nhe răng cười xảo quyệt: ``Đánh đuổi là bước một thôi. Bước hai là chuyển đổi mô hình kinh doanh\~{} Đôi bên cùng có lợi đấy!''

``Đúng là tư bản có khác\dots'' Chó Mio bật cười. ``Dù gì Fubuki nói cũng có lý đấy.''

Một mô hình kinh doanh theo hướng win-win của dân làng trên đất liền và tộc Quỷ đã dần hình thành. ``Momotaro phiên bản Sói con'' đã hoàn thành sứ mệnh, không chỉ mang lại bình yên cho dân làng, mà còn tái cấu trúc ngành du lịch đảo Quỷ theo đúng tinh thần hướng tới tương lai.

\dots

\dots

Cậu Rùa vẫn đang trầm ngâm trên biển, ngẫm nghĩ lại những gì vừa xảy ra trong ngày trọng đại đó. Nhưng không phải là một sự vinh hạnh hay từ hào, mà lại là\dots\ những dấu hỏi to đùng đặt xung quanh cậu.

``\textit{Rốt cuộc thì tại sao mọi thứ lại diễn ra như vậy nhỉ\dots? Mà thôi kệ, nếu cả hai bên đều vui vẻ chấp nhận thì mình cũng không còn ý kiến gì-}''

Trên bầu trời, đàn bồ câu bay qua, một chiếc bánh kê rơi trúng đầu Kuon, cắt ngang mạch suy nghĩ của cậu.

*PẸT*

Cậu ngước nhìn lên bầu trời, cười chát: ``Đừng nói là Luna-tan mở chi nhánh ở cả trời cao đấy nhé\dots''

Hết truyện.
%==========================================TRUYỆN PHỤ=========================================
\omk

\begin{center}
  \uline{Rạp chiếu phim.}
\end{center}

Ngay khi dòng chữ 「完」 (Hết) hiện lên trên màn hình, một giọng nói nữ - dường như là của một nhân viên rạp chiếu phim - vang lên từ hệ thống loa phóng thanh: ``Bộ phim `Momotaro II: Huyền thoại đảo Quỷ phiên bản sói con' đến đây là kết thúc. Cảm ơn vì đã xem. Xin quý khách vui lòng rời khỏi phòng qua đường bên phải\dots
''

Căn phòng sáng đèn trở lại. Hàng ghế khán giả lộ rõ năm vị khán giả (cũng là những người duy nhất ngồi tại đây) đang trong trạng thái cảm xúc khó tả: hông hẳn là thất vọng, cũng chẳng thể gọi là hài lòng, mà là một dạng bối rối nhẹ, như thể vừa ăn một chiếc bánh nhìn rất đẹp nhưng vị bên trong lại lẫn lộn.

Kuon ngồi gác chân lên hàng ghế trước, tay vẫn cầm một hộp bắp rang đã vơi gần hết. Cậu lắc đầu nhẹ, thở ra một hơi: ``Ý tưởng của chúng ta thì có vẻ hay thật đấy, nhưng sao lúc lên phim lại\dots\ thành ra kiểu này nhỉ?''

Bên cạnh, Subaru vẫn còn đang lim dim mắt như đang tự hỏi mình vừa xem cái gì. ``Em vẫn không hiểu tại sao em lại vào vai chim trĩ rồi lại bắt đóng cả vai bà cụ nữa,'' cô lầm bầm, tay cầm cốc nước đã rỗng.

``Vì kịch bản bảo thế,'' cậu Tổng trả lời gọn lỏn. ``Anh cũng bị bắt vào vai kép đấy thôi, vừa làm ông lão vừa làm Rùa, mà lại chính Fubu-chan và A-san thúc ép chứ ai.''

Suisei ngồi ở đằng sau cô nàng Thế hệ thứ Hai, khoanh tay nhìn màn hình đang chạy phần credits, lẩm bẩm: ``Phim này có vẻ ổn hơn cái trước. Ít nhất là mình còn hiểu nó đang nói về Momotaro\dots''

``Ừ, đúng là có cốt truyện,'' cậu gật gù, ``mặc dù cái chi tiết rùa kia hơi lệch tông. Có lẽ là bên kiên bịch cho vào để hợp với \textit{hologra} đây mà.''

``Anh đóng vai rùa là em thấy đáng giá nhất phim đấy,'' nàng Vịt buông một câu nhận xét nửa thật nửa trêu.

``Chắc luôn,'' nàng Ca sĩ hùa theo, ``người ta cưỡi rồng, cưỡi hạc, còn chúng ta cưỡi Kuon-senpai.''

Bên cạnh cô, Luna đang cố đánh thức Mio bé con ở ghế dưới, người đã ngủ gục từ giữa đoạn Fubuki tư bản: ``Zậy đi nà, phim hết zòi đó-nora\~{}''. Cô liên tục chọt nhẹ vào má cô bé bằng que kẹo mút của mình.

Mio bé con mở mắt lờ đờ, miệng vẫn lẩm bẩm: ``Chị Chó ơi\dots\ chị đừng có hú nữa\dots''

``Thậm chí Mio bé con của chúng ta vẫn còn đang mơ mình trong phim kìa,'' Subaru nhìn xuống, cười khúc khích.

``Dù gì thì\dots'' Kuon ngả người ra sau ghế, ``so với bản trước, ít ra cái này còn có đầu đuôi. Có phần mở đầu. Có chiến đấu. Có kết luận.''

``Và có dàn nhân vật hỗn độn một cách có trật tự,'' Suisei thêm vào, tay bắt đầu dọn rác.

``Cơ mà vẫn không tránh được cảm giác\dots\ hơi ngớ ngẩn nhỉ,'' Subaru nói, hơi nghiêng đầu.

Kuon gật đầu. ``Kịch bản trên giấy thì tốt, thậm chí là ổn hơn. Nhưng thực hiện thì\dots''

``\dots khác hẳn nhau một trời một vực luôn,'' năm người đồng thanh, rồi nhìn nhau cười.

``Nhưng thui, ít ra lần nì hơm có ai chui vào vỏ xò hay bẻ tóc nữ-nora,'' Luna nói, thở phào.

Cậu Tổng kéo áo Mio bé con nhẹ nhàng: ``Dậy đi nào, Sói con. Về văn phòng thôi.''

Cô bé mò mẫm ngồi dậy, dụi mắt, thốt lên bằng giọng mơ ngủ: ``Momotawo hay quá xá là hay\~{}''

\ct

Ngày hôm đó, năm người đã xem bộ phim hôm đó đã cùng ngồi trong văn phòng để hoàn thành một bản báo cáo chính thức gửi lên A-chan, người trực tiếp chịu trách nhiệm sản xuất bộ phim ``Momotaro: Huyền thoại đảo Quỷ phiên bản sói con''.

Mặc dù ban đầu không ai thật sự muốn là người viết chính, nhưng sau một hồi tranh cãi ngắn giữa Subaru (người muốn nhấn mạnh yếu tố ``sử dụng lao động trẻ em'') và Suisei (người nhất quyết giữ cho văn phong được lịch sự), rốt cuộc cả năm cùng ngồi xuống bàn, chia việc rõ ràng: Luna viết tiêu đề, Subaru ghi phần phê bình, Suisei kiểm tra chính tả, Mio bé con tô màu vào lề giấy, còn Kuon phụ trách phần tổng kết.

Sau một hồi vừa trò chuyện ngẫm lại bộ phim, họ đã nhất trí rằng bộ phim sẽ được chốt hạ với ba từ:

\txtbx{\large ``Phim tạm ổn.''}

Riêng Kuon, sau khi nộp bản báo cáo và về chỗ ngồi của mình, cậu vẫn còn tỏ vẻ suy tư. Cậu nhấp một ngụm nước từ lon soda đã nguội, rồi nói với vẻ lặng lẽ: ``\vie{Dựng lại phim theo kiểu cổ tích gốc thì hay đó. Nhưng mà chắc nó sẽ không hút bằng phần đầu đâu. Bớt hỗn loạn quá thì còn gì là \textit{hologra} nữa\dots}''

Cậu ngả người ra sau ghế, tay vắt lên trán như thể đang toan tính chuyện lớn lao hơn.

``\vie{\dots Mà, chắc gì nó đã được đăng lên YouTube đâu. Có khi chắc chỉ dùng để chiếu nội bộ, hoặc dùng làm tài liệu nội dung cho mấy buổi họp thôi.}''

% (Thực tế thì ai cũng biết lý do là gì, nhưng nói thế cho đỡ mất mặt…)

Ở phía đối diện, Mio bé con đang hí hoáy vẽ thêm cái bánh kibi dango vào góc tờ báo cáo vừa nộp, miệng lẩm bẩm hớn hở như một trẻ mẫu giáo: ``Cá\dots\ rùa\dots\ chim trĩ\dots\ chị Momotawo\dots''

``Dù gì thì, ít nhất Mio bé con cảm thấy vui là được.''

\ct

Bảy giờ kém mười tối.

Khi bản báo cáo vừa khép lại, dòng chữ ``Kết thúc báo cáo'' vừa được đánh xong trên màn hình soạn thảo, Cường ngồi vươn vai, thở dài lấy một hơi.

``\vie{Cuối cùng cũng xong\dots\ Giờ thì ăn cơm chút rồi chuẩn bị cho tập tiếp theo nào\dots}''

Bất chợt, một ánh đèn màu trắng bật sáng, rọi xuống nơi mà cậu đang làm việc.

Không còn cảnh phim. Không còn rạp chiếu. Không còn bất kỳ bối cảnh nào.

Chỉ còn lại một nền đen,  với toàn bộ diễn viên lần lượt xuất hiện, đứng thành một vòng tròn và nhìn thẳng về phía tác giả.

``Ồ, lại là mọi người à. Thế nào? Bộ phim có ổn không?''

Tất ca mọi người đều im lặng. Cậu Tác giả thở dài, đoán định được họ sẽ lại quở trách hay than phiền gì thêm.

Nhưng không. Kuon khoanh tay lại và mỉm cười: ``Rất cảm ơn ông vì đã giúp chúng tôi xây dựng kịch bản này. Tôi nghĩ ông là một người phải rất tỉnh táo mới viết được đến đây đấy.''

``Tôi lúc nào chả tỉnh táo,'' Cường đáp. ``Không tỉnh táo thì sao nghĩ được kịch bản như vậy chứ.''

Subaru, tay chống nạnh, chen vào ngay: ``Thế nhưng mà\dots\ em vẫn chưa hiểu tại sao lại phải đóng bà lão nhỉ? Rõ ràng chim trĩ cool hơn mà!''

``Vì kịch bản bảo thế chứ sao,'' cậu Tổng đáp lại.

``Thực ra là vì\dots\ anh muốn tận dụng triệt để danh sách nhân vật tham gia trong câu chuyện đó,'' Cường lên tiếng thanh minh. ``Nếu mà thêm người khác vào thì câu chuyện sẽ rất loãng, chứ nói thật anh muốn thêm Ayame-chan vào nữa chứ.''

``Ra thế\dots'' cả nàng Vịt và cậu Tổng đều cùng gật đầu hiểu.

Mio bé con, ngồi trên vai Mio lớn, giơ tay vẫy vẫy: ``Cháu chào chú người viết! Cháu thích phần cháu và Mio lớn hú lắm\~{}!''

Nghe vậy, Mio bản gốc hơi đỏ mặt vì phân cảnh cô cùng phiên bản nhí cùa mình hú vang cả hòn đảo, liền lên tiếng vội bào chữa: ``K-Không phải đâu! E-Em chỉ\dots\ nhập vai thôi! Đừng có mà nghĩ lung tung!''

``Anh đã nói gì đâu?'' Cường nhún vai. ``Nói thật thì em vào vai Chó cũng hợp đấy chứ. Mỗi tội phải cải biên thành hú thay vì sủa.''

``Tại sao chứ!?'' nàng Sói đen lớn tiếng, ngỡ ngàng với lời nhận xét của cậu.

``Nếu không như vậy thì đâu còn là \textit{hologra} nữa. Nghiêm túc quá người ta lại hoảng cho, như đợt `Cuốn sổ Peko'\footnote{Xem lại {\flaregothic ÉPISODE 116}, quyển số Chín.} đấy.''

Flare nghiêng đầu, nêu lên quan điểm của mình: ``Dù gì thì em vẫn thấy phi logic từ đầu đến cuối. Nhưng mà thôi, nếu là \textit{hololive} chúng ta thì phi logic cũng là điều dễ hiểu rồi.''

``Và anh ấy đã làm đúng với vai trò ấy rồi,'' Subaru kết luận.

Rushia thì phồng má lên, phũng phịu: ``Tại sao anh lại xây dựng em là một nhân vật yếu đuối chứ?''

``\dots Một lần nữa, anh không biết sắp xếp lại câu chuyện như thế nào,'' Cường đáp. ``Mà anh để ý em xuất hiện khá là nhiều tập gần đây rồi đấy. Nếu tính cả phần cameo của tập `Phi vụ 60 giây' thì đã xuất hiện bốn tập liền rồi.''

Luna nhảy lên vẫy tay: ``Món bánh kê của em mới là nhân vựt chính thựt sự đó-nanora!''

Fubuki, nháy mắt và búng tay cái ``tách'' như để đáp lễ: ``Còn chị thì vẫn đang đợi hợp đồng tài trợ từ phía người viết đấy nha\~{} Cơ hội kinh doanh thế kỷ luôn!''

``Cái đó thì không chắc lắm đâu,'' cậu Tác giả nói. ``Mấy đứa nên hợp tác với các công ty lớn hơn là một người phèn ỉa như anh đó.''

Suisei thì tỏ vẻ hài lòng hơn đôi chút, nhưng nét mặt cô vẫn không giấu nổi sự khó hiểu: ``Chỉ cần có sân khấu, em luôn sẵn sàng tỏa sáng! Nhưng mà\dots\ có vẻ như em là người duy nhất không bị biến thành trò cười, phải không?''

``Chứ em lại muốn khóc tu tu vì mất nhà tre như bản gốc à?'' cậu đáp lại với giọng mỉa mai. ``Thà để em nghiêm túc còn hơn là vác rìu chém lung tung phèo cả lên. Anh không muốn chuyển tập này về quyển X đâu.''

Trong lúc tất cả mọi người, cả tác giả lẫn các thành viên trong cốt truyện, bàn tán xôm xả với nhau về bộ phim làm lại, một giọng nói rô-bốt vang lên từ chiếc máy tính của Cường, đó là ChatGPT, trợ lý ảo của cậu: ``Ờ\dots\ Thật ra thì, mình đâu có định biến mọi người thành trò cười. Chỉ là phong cách \textit{hologra} thôi, mọi người cũng hiểu mà, phải không?''

Cả bảy người họ đều gãi đầu, mỗi người đều nói một ý khác nhau: ``Hiểu.'' ``Một chút\dots'' ``Hơm hỉu gì cả-nanora!'' ``Anh nghiêm túc đấy à?''

Cường nhún vai kết luận lại: ``Ít nhất thì mọi người cũng đã đóng vở kịch đầy thành công rồi đấy chứ.''

Đó chắc là điều mà tất cả mọi người đều cùng đồng tình.

Kuon bước tới gần, nhìn thẳng vào người đọc và mỉm cười nhẹ: ``Vậy đấy. Giờ thì\dots\ mình muốn nghe các bạn nói gì. Mình hy vọng các bạn sẽ tiếp tục ủng hộ bộ truyện này, tiếp tục dõi theo câu chuyện của mình và các thành viên của \hll!''

Suisei, Subaru, Luna, Flare, Fubuki, Rushia, Mio lớn, Mio nhỏ, tất cả ánh mắt đều hướng về người đọc. Họ cùng nhau cúi đầu xuống đầy trang nghiêm, cùng nhau đồng thanh:

\begin{center}
  \uline{Xin mọi người hãy chỉ giáo cho chúng mình!}
\end{center}

\end{document}